\documentclass[11pt]{article}
\usepackage[margin=1.1in]{geometry}
\usepackage{amsmath,amssymb,amsthm}
\usepackage{booktabs}
\usepackage[round,authoryear]{natbib}
\usepackage[hidelinks]{hyperref}
\usepackage{microtype}
\makeatletter
\renewcommand\paragraph{\@startsection{paragraph}{4}{\z@}%
  {3.25ex \@plus1ex \@minus.2ex}{-1em}%
  {\normalfont\normalsize\itshape}}
\makeatother
\newtheorem{proposition}{Proposition}
\theoremstyle{remark}
\newtheorem{remark}{Remark}
\newcommand{\CV}{\mathrm{CV}}

\title{The Immunity Factor from Contact Surveys:\\ Where Epidemics Turn Over}
\author{Peter Cotton\thanks{\texttt{peter.cotton@microprediction.com}. Code, data sources and a dated log of
results: \texttt{github.com/microprediction/epidemiology}; summaries at
\texttt{epidemiology.microprediction.org}.}}
\date{September 2026}

\begin{document}
\maketitle

\begin{abstract}
Many COVID-19 waves reached their peak of incidence when estimated cumulative infections were well below the
threshold $1 - 1/R$ implied by their initial growth under a homogeneous model. Because people with many contacts are
infected sooner, the effective reproduction number falls faster than the susceptible share $S$, approximately as
$R\,S^{\lambda}$, where the immunity factor $\lambda$ is one in a homogeneous population; the peak then comes at about
$1 - R^{-1/\lambda}$. In repeated UK contact surveys, a model of persistent and changing activity predicts
$\lambda = 4.6$--$5.2$ when daily contacts are capped at 50, the survey's own convention, and 3.4--3.7 at a cap of 20,
without fitting to epidemic outcomes. Winter 2020--21 US county waves imply 2.4--5.3 for infection fatality rates from
0.5\% to 1\%, and about 3.5 at 0.7\%.
\end{abstract}

\section{Introduction}

In 341 US county epidemic waves of winter 2020--21, the median wave reached its peak of incidence when estimated
cumulative infections were 16\% of the population. Under a homogeneous model, the growth at the start of those waves
implies a median threshold of 30\%. Infections here are estimated from deaths at a fatality rate of 0.7\%, and growth
is converted to a reproduction number with a generation interval of mean 5.5 days. We call the peak of incidence the
\emph{turnover}.

In the spring of 2020 case counts doubled in under a week, which a homogeneous model reads as a reproduction number
near 2.5 and a threshold of 60\%. \citet{cotton2020fundamental,cotton2020herd} called the gap between such thresholds
and the observed peaks the herd-immunity paradox and argued that it follows from variation in transmission across people,
places and time: early growth is dominated by the fastest transmitters, and the turnover by a harmonic mean
\citep{cotton2020herd}.

People with many contacts are infected sooner, and while infectious they generate more exposure, so each infection
removes more than its share of future transmission. \citet{britton2020} and \citet{gomes2022} made this quantitative
for persistent differences in activity or susceptibility. \citet{tkachenko2021} added activity that changes over time,
named the immunity factor, and showed that the resulting suppression is partly transient; \citet{tkachenko2021elife}
modelled each person's activity as a mean-reverting process.

The effect is summarized by the \emph{immunity factor} $\lambda$, the rate at which the effective reproduction number
falls relative to the susceptible share. We estimate it in two ways that share no data. The first uses repeated
contact surveys and a model of persistent and changing activity, without fitting to epidemic outcomes. The second
backs a value out of the peaks of US county waves under an assumed fatality rate. The two ranges overlap, and both
depend on their assumptions: the contact cap on one side, the fatality rate and the growth proxy on the other.
\citet{gomes2022} obtain an independent value near 3.5 by fitting heterogeneous connectivity to deaths in England and
Scotland.

Where contacts were cut by lockdowns or voluntary distancing, peaks came with 1--5\% infected. The value of $\lambda$
implied by such a peak absorbs the drop in contacts and does not measure immunity.

\section{The turnover rule}

\paragraph{Unequal activity.} The immunity factor and the relations \eqref{eq:power}--\eqref{eq:rule} below are due
to \citet[eqs.~9, 14, 15]{tkachenko2021}. Let $S$ be the susceptible share. In a homogeneous population the effective reproduction
number is $R_e = R\,S$. For fixed gamma-distributed activity the power law below is exact, as shown next. In other
models we define $\lambda$ by the initial fall of $R_e$ and use the power law as an approximation to predict the
turnover. When people differ in activity, write
\begin{equation}\label{eq:power}
  R_e = R\,S^{\lambda}.
\end{equation}
Early on, $R_e \approx R\,(1 - \lambda\,(1 - S))$, so each percentage point of infection lowers $R_e$ by $\lambda$
percentage points of $R$. Setting $R_e = 1$ gives the infected share at the turnover,
\begin{equation}\label{eq:rule}
  a^\star = 1 - R^{-1/\lambda}.
\end{equation}
The textbook threshold is the case $\lambda = 1$. Equivalently, the textbook formula applies with the adjusted
reproduction number $R^{1/\lambda}$.

Here $R$ is the initial effective reproduction number of the model. When $R$ is inferred from growth it also depends
on the assumed generation interval, and in a model with changing activity the two need not coincide.

\paragraph{Gamma activity.} Let activity $a$ have a gamma distribution with mean one and squared coefficient of
variation $\CV^2 = 1/k$, and let it scale both the chance of infection and the number of people infected. After
cumulative force of infection $Z$, the susceptible share is $S = (1 + Z/k)^{-k}$ and $R_e \propto \mathbb{E}[a^2
e^{-aZ}] \propto (1 + Z/k)^{-k-2}$. Hence \eqref{eq:power} holds exactly with
\begin{equation}\label{eq:gamma}
  \lambda = 1 + 2\,\CV^2 ,
\end{equation}
as in \citet[eqs.~21--22]{gomes2022}, who give $1 + \CV^2$ when only susceptibility varies, and
\citet[eq.~14]{tkachenko2021}.

\paragraph{Age structure.} Let $M_{ij}$ be the mean number of daily contacts of a person in age group $i$ with people
in group $j$, and $n_i$ the population shares. Infected fractions grow along the leading right eigenvector $v$ of $M$;
let $u$ be the left eigenvector. Depleting susceptibles along $v$ lowers the leading eigenvalue at the rate
\begin{equation}\label{eq:lamage}
  \lambda_{\text{age}} = \frac{\sum_i u_i v_i^2 \big/ \sum_i u_i v_i}{\sum_i n_i v_i}.
\end{equation}
The numerator measures how strongly early infections sit in the groups that drive transmission; the denominator
converts to the overall infected share. The same formula applies to any finite type structure, and we use it for age
crossed with activity.

\paragraph{Accuracy of the rule.} Equation \eqref{eq:lamage} is a first-order slope, so \eqref{eq:rule} is an
approximation away from the gamma case. In the joint age and activity models of Section~\ref{sec:contacts}, solved
numerically to the peak, the rule understates the infected share at the turnover by 0.6 to 1.2 percentage points at
$R = 1.3$, 1.4 to 2.9 points at $R = 1.6$, and 3.6 to 7.0 points at $R = 2.5$.

\section{Activity that changes over time}\label{sec:switching}

\paragraph{A solvable model.} \citet{tkachenko2021elife} let each person's activity follow a mean-reverting
square-root process with a correlation time. A two-state version has closed forms. Let each person's activity take two values with long-run shares $p_1, p_2$, mean one and
variance $\sigma^2$, so that $(a - 1)^2 = c\,(a - 1) + \sigma^2$ for a constant $c$ that sets the skew. Activity
switches on a Markov clock with generator $Q$ and relaxes at rate $\kappa$. A person with activity $a$ is infected at
rate $\beta\,a\,\Phi$, where $\Phi = \sum_k a_k I_k$, and infectious people recover at rate $\gamma$.

\begin{proposition}[Growth, reproduction number and immunity factor]\label{prop:switch}
The early growth rate $r$ and the basic reproduction number solve
\begin{equation}\label{eq:growth}
  1 = \beta\Big(\frac{1}{r + \gamma} + \frac{\sigma^2}{r + \gamma + \kappa}\Big), \qquad
  R_0 = \beta\Big(\frac1\gamma + \frac{\sigma^2}{\gamma + \kappa}\Big).
\end{equation}
During growth, with $\rho = r/(r + \kappa)$, the immunity factor is
\begin{equation}\label{eq:lamswitch}
  \lambda = \frac{(1 + \sigma^2\rho)/\gamma + \sigma^2\,(1 + (1 + c)\rho)/(\gamma + \kappa)}{1/\gamma + \sigma^2/(\gamma + \kappa)} .
\end{equation}
\end{proposition}

The proof is in Appendix~\ref{app:proof}. The resolvent $(r + \gamma - Q^\top)^{-1}$ acts on only two directions,
the stationary law $p$ and the deviation $(a - 1)\circ p$, which decays at rate $\kappa$. People with more contacts
are infected sooner and, while infectious, generate more exposure. The second term in \eqref{eq:growth} is the
persistence of an individual's activity over the infectious period: it contributes while a busy phase outlasts the
infection.

\paragraph{Limits.} With $\kappa = 0$, \eqref{eq:lamswitch} is $\langle a^3\rangle/\langle a^2\rangle$, the value for
persistent heterogeneity \citep[eqs.~12, 18]{tkachenko2021}. As $\kappa \to \infty$ it tends to one. For fast switching,
\begin{equation}\label{eq:gk}
  \lambda - 1 \approx r\,K_a, \qquad K_a = \frac{\sigma^2}{\kappa},
\end{equation}
where $K_a$ is the Green--Kubo number of activity, the integral of its autocovariance. A fast epidemic outruns the
reshuffling of activity and uses up the busy phase before it is refilled.

\paragraph{Effective variance.} For the prediction in Section~\ref{sec:contacts} we combine a persistent variance
$\sigma_p^2$ with a switching variance $\sigma_s^2$ of correlation time $\tau$ through
\begin{equation}\label{eq:cveff}
  \CV^2_{\text{eff}} = \sigma_p^2 + \frac{r}{r + 1/\tau}\,\sigma_s^2 ,
\end{equation}
the persistent part acting in full and the switching part in proportion to the share of a busy phase that survives
one growth time. This interpolates between the limits of Proposition~\ref{prop:switch} and is an approximation.

\section{Mixtures of places}

\citet{cotton2020herd} reaches the same conclusion from a population made of many places. To isolate the effect of
heterogeneous local rates, consider independent places with equally seeded epidemics, each with transmission rate
$\beta = \bar\beta\,\eta$ where $\eta$ has mean one across places. Early growth of the total reads a transmission
rate $\beta_g = t^{-1}\log\mathbb{E}\,e^{t\bar\beta\eta} = G\,\bar\beta$ with $G > 1$. The susceptible share at the
turnover, averaged over places, is \citep[\S\S3.1--3.2]{cotton2020herd} $\mathbb{E}[\gamma/\beta] = J/\bar R$ with $J = \mathbb{E}[1/\eta] \ge 1$. Together,
\begin{equation}
  s^\star \approx \frac{J\,G}{R_g},
\end{equation}
larger than the textbook $1/R_g$ by two convexity adjustments. In \citet{cotton2020herd} the rate follows an
Ornstein--Uhlenbeck process, its growth convexity is Vasicek's bond-price formula, and the long-run growth excess is
the variance times the correlation time \citep[eq.~16]{cotton2020herd}. The same holds for a two-state switch. When
each place's rate switches between $\bar\beta \pm
\delta$ at rate $\nu$, the total grows at the top eigenvalue $\bar\beta - \gamma - \nu + \sqrt{\nu^2 + \delta^2}$, which
is $\bar\beta - \gamma + \delta^2/(2\nu)$ to first order: the growth convexity is again a Green--Kubo number.

\section{The immunity factor from contact surveys}\label{sec:contacts}

\paragraph{One-day diaries.} The POLYMOD diaries \citep{mossong2008} record the contacts of 7,290 people in eight
European countries on one day. Diary space limited entries to 29 contacts (30 in the United Kingdom), and in four
countries respondents estimated their professional contacts rather than listing them. We count listed contacts, which
understates the upper tail. Within age groups the squared coefficient of variation of the number of contacts is
0.53 for ages 0--19, 0.63 for 20--64 and 0.77 for 65 and over. Combining these with the age matrices of
\citet{prem2021} through \eqref{eq:lamage}, with activity held fixed, gives $\lambda = 2.70$ for the United States,
2.77 for Sweden, 2.66 for Spain, 3.02 for the United Kingdom and 2.30 for Brazil. Age alone gives 1.3--1.7.

A one-day diary cannot tell a busy person from a busy day, so these values mix lasting and passing variation.
\citet{gomes2022} estimate the coefficient of variation from POLYMOD and nineteen other surveys, make the same
point, and use two-day diaries to separate the two.

\paragraph{Repeated surveys.} The CoMix survey \citep{jarvis2020,gimma2022,gimmadata2024} ran in the United Kingdom
from March 2020 to March 2022. Each participant reported their contacts up to eight or ten times, two weeks apart. We
use 21,158 adults aged 18--69 seen at least three times, 160,869 survey days in all. The median participant is
observed over 123 days and 90\% within 267 days; about 4\% rejoin after a gap of around seven months.

Activity is the number of contacts on the diary day divided by the mean for the same week and age group. This removes
the common effect of lockdowns and seasons on each age group, though not effects that differ between people. The
covariance of a person's activity between surveys $\ell$ days apart is fitted as
\begin{equation}
  C(\ell) = \sigma_p^2 + \sigma_s^2\,e^{-\ell/\tau} \quad (\ell \ge 7), \qquad C(0) = \sigma_p^2 + \sigma_s^2 + \sigma_n^2 ,
\end{equation}
with a persistent part, a switching part of correlation time $\tau$, and a residual $\sigma_n^2$ that repeat
observations at lags of a week or more cannot resolve.

\begin{table}[h]
\centering
\begin{tabular}{rrrrrr}
\toprule
contacts capped at & one-day $\CV^2$ & persistent $\sigma_p^2$ & switching $\sigma_s^2$ & $\tau$ (days) & residual $\sigma_n^2$ \\
\midrule
20 & 1.82 & 0.54 (30\%) & 0.23 (13\%) & 36 & 1.05 (57\%) \\
50 & 3.90 & 1.06 (27\%) & 0.51 (13\%) & 28 & 2.34 (60\%) \\
100 & 6.54 & 1.69 (26\%) & 0.88 (14\%) & 23 & 3.97 (61\%) \\
\bottomrule
\end{tabular}
\caption{Decomposition of adult activity in the CoMix UK panel.}\label{tab:comix}
\end{table}

Under this covariance model, 26--30\% of the variance of capped daily contacts is persistent and 13--14\% belongs to a
component with a correlation time of 23 to 36 days. The remaining 57--61\% is unresolved by repeat observations at lags
of a week or more. It includes day-to-day variation, and it may also include genuine correlations lasting less than a
week, which could matter over an infectious period. The shares are similar under the three caps, but the levels are
not. The fitted covariance is flat from one month to the longest lags, so the persistent component behaves like a
habit over the months each person is observed. Restricting the fit to lags of at most 200 days changes each share by
at most one percentage point.

\paragraph{Prediction.} Inserting \eqref{eq:cveff} into the joint age and activity structure with the UK matrix gives
the immunity factors in Table~\ref{tab:predict}, for growth rates from 0.02 to 0.2 per day. The residual is left
out, so these values treat all sub-weekly variation as irrelevant to depletion. No epidemic outcome enters.

\begin{table}[h]
\centering
\begin{tabular}{rccc}
\toprule
contacts capped at & $\lambda$ & turnover at $R = 1.6$ & turnover at $R = 2.5$ \\
\midrule
20 & 3.4--3.7 & 12--13\% & 22--24\% \\
50 & 4.6--5.2 & 9--10\% & 16--18\% \\
100 & 5.9--6.8 & 7--8\% & 13--14\% \\
\bottomrule
\end{tabular}
\caption{Immunity factor predicted from the CoMix decomposition and the UK age matrix, without fitting to epidemic outcomes.}\label{tab:predict}
\end{table}

The immunity factor changes by only about 7\% across a tenfold range of growth rates, because the switching part is
small. The level depends on how the few very high contact counts are treated: it rises from 3.4--3.7 at a cap of 20 to
5.9--6.8 at 100. The CoMix team caps at 50 contacts a day, the value they found to track changes in the reproduction
number most closely \citep{gimma2022}. How many of a very busy person's contacts can transmit is unresolved.

\section{The immunity factor in epidemics}

\paragraph{US county waves.} We identify 471 epidemic waves in 372 US counties of more than 50,000 people from daily
deaths \citep{nyt2021}, February 2020 to June 2021. The data count deaths by the date
they were reported, which adds noise to the timing of death peaks. For each wave, early growth $r$ gives $R$ through
the relation $R = 1/M(-r)$ of \citet{wallinga2007}, with $M$ the moment generating function of an assumed gamma
generation interval of mean 5.5 days and standard deviation 2.1 days. Cumulative deaths divided by an
infection fatality rate give the infected share when the wave began, $a_0$, and at the death peak, $a_1$.

Suppose contacts stayed at their early-wave level. Then the wave's growth and its infected share at the peak imply,
through \eqref{eq:rule}, the value
\begin{equation}
  \hat\lambda = \frac{\ln R}{-\ln\big((1 - a_1)/(1 - a_0)\big)} ,
\end{equation}
which we call the \emph{implied peak index}. It equals the immunity factor only if transmission per contact and the
number of contacts did not change during the wave. When contacts fell, $\hat\lambda$ absorbs the fall and does not
measure immune selection.

Between 95\% and 100\% of waves in every season turned over below the homogeneous threshold. In winter 2020--21,
with 341 waves and growth estimated from deaths, a constant $\hat\lambda$ fits at 3.5 for a fatality rate of 0.7\%,
2.4 at 0.5\% and 5.3 at 1\%. With growth estimated from cases it is about 2.1 at 0.7\%.

\paragraph{The slope is not identified.} Grouped by growth rate, the winter waves infected nearly the same share of
their susceptibles whatever their $R$. Growth read from deaths and from cases three weeks earlier correlate at only
0.23 in winter, and noise in a regressor flattens any fitted slope. With case growth as an instrument the slope of
log depletion on $\log\ln R$ is 0.91 in spring (90\% interval 0.46 to 1.50) and 2.9 in winter (0.7 to 12.5). The
county data identify the level of $\hat\lambda$ but not its dependence on growth, which Table~\ref{tab:predict} says is
small.

\paragraph{Serology.} Antibody surveys measure the infected share directly. Scaling each survey back to the death peak
with cumulative deaths gives the infected share at the turnover and the implied peak index in Table~\ref{tab:sero}.

\begin{table}[h]
\centering
\small
\begin{tabular}{lllccc}
\toprule
place & survey & contacts & $R$ & infected & $\hat\lambda$ \\
\midrule
England, 9 regions & \citet{ward2021} & lockdown & 2.1--3.5 & 0.9--5.0\% & 17--94 \\
Geneva & \citet{stringhini2020} & closures & 3.3 & 5.1\% & 23 \\
Stockholm & \citet{castrodopico2021} & voluntary & 2.8 & 2.3--3.1\% & 33--44 \\
Spain, spring & \citet{pollan2020} & lockdown & 1.9--4.3 & 0.3--4.6\% & about 60 \\
Spain, autumn & \citet{sanidad2020} & curfew & 1.1--1.5 & 2.6\%$^\dagger$ & about 7 \\
Manaus & \citet{buss2021} & limited & 1.6 & about 17\% & 2.5 \\
\bottomrule
\end{tabular}
\caption{Turnovers anchored by serology; the infected share is at the turnover. $^\dagger$Share of those still susceptible after the first wave.}\label{tab:sero}
\end{table}

Where contacts were cut, by law or voluntarily, the first wave turned over with 1--5\% infected and $\hat\lambda$
exceeds 15. These values are consistent with the observed reduction in contacts and cannot be read as immunity. Where
control was weaker, in Manaus and in the US county waves of winter 2020--21, $\hat\lambda$ is 2 to 3.5.
\citet{tkachenko2021} fit a transient immunity factor of about 4 to New York City and Chicago during stay-at-home
orders, with a model that includes mitigation. \citet{buss2021} report that distancing in Manaus worked together with
immunity.

\section{Discussion}

Contact surveys predict an immunity factor of 4.6--5.2 at the survey's own cap of 50 contacts a day, 3.4--3.7 at a cap
of 20 and 5.9--6.8 at 100. Winter county waves imply a peak index of 2.4--5.3 over fatality rates from 0.5\% to 1\%,
3.5 at 0.7\%, and about 2.1 with growth from cases. \citet{gomes2022} fit about 3.5 to England and Scotland. The
ranges overlap. At $R = 2.5$, $\lambda = 3.5$ puts the turnover near 23\% and $\lambda = 5$ near 17\%.

On the contact side the level depends on how very high contact counts are treated and on the residual below one week.
On the epidemic side it scales with the assumed fatality rate and the growth proxy. The two sets of uncertainties are
of similar size.

In the CoMix data most of the resolved heterogeneity persists over the months each person is observed. The switching
part is smaller and decorrelates over about a month, so in this model a first-wave turnover is mostly a lasting
reduction of the threshold with a transient on top.

\citet{tkachenko2021} reach a different balance. They put the long-run immunity factor near 2 and conclude that most of
the first-wave suppression was transient. A survey that follows the same people for years would settle how long the
persistent component lasts.

The reported later attack in Manaus, 66--76\%, is difficult to reconcile with a fixed-activity model of this size.
Changing activity is one possible explanation. The contact distribution in Manaus and the blood-donor estimates are
additional uncertainties.

The prediction can be tested out of sample. CoMix surveys ran in 21 European countries \citep{wong2023}, though adult
participants in most of them reported only seven times. Fixing the treatment of high contact counts in advance, each
survey predicts an immunity factor that the country's serology and death curves can confirm or refute.

\appendix
\section{Proof of Proposition~\ref{prop:switch}}\label{app:proof}

Let $I = (I_1, I_2)$ be the infectious people by current activity and $\Phi = a^\top I$. Early on, $\dot I = \beta\,\Phi\,(a
\circ p) - \gamma I + Q^\top I$. Seek $I = v\,e^{rt}$:
\begin{equation}
  v = \beta\,\Phi\,(r + \gamma - Q^\top)^{-1}(a\circ p).
\end{equation}
Since $Q^\top p = 0$ and $Q^\top\big((a - 1)\circ p\big) = -\kappa\,(a - 1)\circ p$ for a two-state chain, and $a\circ p = p +
(a - 1)\circ p$,
\begin{equation}
  (r + \gamma - Q^\top)^{-1}(a\circ p) = \frac{p}{r + \gamma} + \frac{(a - 1)\circ p}{r + \gamma + \kappa}.
\end{equation}
Taking $a^\top$ of both sides, with $a^\top p = 1$ and $a^\top\big((a - 1)\circ p\big) = \sigma^2$, gives \eqref{eq:growth}; $R_0$
is the case $r = 0$.

For the immunity factor, the susceptible deficit $s = p - S$ solves $\dot s = \beta\,\Phi\,(a \circ S) + Q^\top s$, and
during growth $s = \beta\,\Phi\,\big(p/r + (a - 1)\circ p/(r + \kappa)\big)$ to first order, with overall attack $x = \beta\Phi/r$.
Then $a\circ s$ has components $x\,(1 + \sigma^2\rho)$ along $p$ and $x\,(1 + (1 + c)\rho)$ along $(a - 1)\circ p$, using
$(a - 1)^2 = c\,(a - 1) + \sigma^2$. The next-generation number with susceptibles $S$ is $\beta\,a^\top(\gamma - Q^\top)^{-1}(a\circ S)$,
and subtracting its value at $S = p$ gives $R_e = R_0\,(1 - \lambda x)$ with $\lambda$ as in \eqref{eq:lamswitch}.
The formulas agree with numerical solution of the model: the growth rate to four digits and $\lambda$ to 0.1--0.3\%
(\texttt{04\_activity\_switching.py}).

\begin{thebibliography}{99}
\bibitem[Britton et~al.(2020)]{britton2020} Britton, T., Ball, F. and Trapman, P. (2020). A mathematical model reveals the influence of population heterogeneity on herd immunity to SARS-CoV-2. \emph{Science} 369(6505):846--849.
\bibitem[Buss et~al.(2021)]{buss2021} Buss, L. F., Prete, C. A., Abrahim, C. M. M. et al. (2021). Three-quarters attack rate of SARS-CoV-2 in the Brazilian Amazon during a largely unmitigated epidemic. \emph{Science} 371(6526):288--292.
\bibitem[Castro Dopico et~al.(2021)]{castrodopico2021} Castro Dopico, X., Muschiol, S., Christian, M. et al. (2021). Seropositivity in blood donors and pregnant women during the first year of SARS-CoV-2 transmission in Stockholm, Sweden. \emph{Journal of Internal Medicine} 290(3):666--676.
\bibitem[Cotton(2020a)]{cotton2020fundamental} Cotton, P. (2020a). A Fundamental Theorem for Epidemiology. LinkedIn, 27 May 2020.
\bibitem[Cotton(2020b)]{cotton2020herd} Cotton, P. (2020b). Addressing the herd immunity paradox using symmetry, convexity adjustments and bond prices. arXiv:2006.07341.
\bibitem[Gomes et~al.(2022)]{gomes2022} Gomes, M. G. M., Ferreira, M. U., Corder, R. M., King, J. G. et al. (2022). Individual variation in susceptibility or exposure to SARS-CoV-2 lowers the herd immunity threshold. \emph{Journal of Theoretical Biology} 540:111063.
\bibitem[Gimma et~al.(2022)]{gimma2022} Gimma, A., Munday, J. D., Wong, K. L. M. et al. (2022). Changes in social contacts in England during the COVID-19 pandemic between March 2020 and March 2021 as measured by the CoMix survey: a repeated cross-sectional study. \emph{PLOS Medicine} 19(3):e1003907.
\bibitem[Gimma et~al.(2024)]{gimmadata2024} Gimma, A., Wong, K. L. M., Van Zandvoort, K., Jarvis, C. I. and Edmunds, W. J. (2024). CoMix social contact data (UK). Zenodo, doi:10.5281/zenodo.13684044.
\bibitem[Jarvis et~al.(2020)]{jarvis2020} Jarvis, C. I., Van Zandvoort, K., Gimma, A. et al. (2020). Quantifying the impact of physical distance measures on the transmission of COVID-19 in the UK. \emph{BMC Medicine} 18:124.
\bibitem[Mossong et~al.(2008)]{mossong2008} Mossong, J., Hens, N., Jit, M. et al. (2008). Social contacts and mixing patterns relevant to the spread of infectious diseases. \emph{PLoS Medicine} 5(3):e74.
\bibitem[New York Times(2023)]{nyt2021} The New York Times (2023). Coronavirus (Covid-19) data in the United States. \texttt{github.com/nytimes/covid-19-data}, archived 24 March 2023, accessed September 2026.
\bibitem[Ministerio de Sanidad(2020)]{sanidad2020} Ministerio de Sanidad and Instituto de Salud Carlos III (2020). Estudio ENE-COVID: cuarta ronda. Estudio nacional de sero-epidemiolog\'ia de la infecci\'on por SARS-CoV-2 en Espa\~na. Informe, 15 December 2020.
\bibitem[Poll\'an et~al.(2020)]{pollan2020} Poll\'an, M., P\'erez-G\'omez, B., Pastor-Barriuso, R. et al. (2020). Prevalence of SARS-CoV-2 in Spain (ENE-COVID): a nationwide, population-based seroepidemiological study. \emph{The Lancet} 396(10250):535--544.
\bibitem[Prem et~al.(2021)]{prem2021} Prem, K., van Zandvoort, K., Klepac, P. et al. (2021). Projecting contact matrices in 177 geographical regions: an update and comparison with empirical data for the COVID-19 era. \emph{PLOS Computational Biology} 17(7):e1009098.
\bibitem[Stringhini et~al.(2020)]{stringhini2020} Stringhini, S., Wisniak, A., Piumatti, G. et al. (2020). Seroprevalence of anti-SARS-CoV-2 IgG antibodies in Geneva, Switzerland (SEROCoV-POP): a population-based study. \emph{The Lancet} 396(10247):313--319.
\bibitem[Tkachenko et~al.(2021a)]{tkachenko2021} Tkachenko, A. V., Maslov, S., Elbanna, A., Wong, G. N., Weiner, Z. J. and Goldenfeld, N. (2021a). Time-dependent heterogeneity leads to transient suppression of the COVID-19 epidemic, not herd immunity. \emph{Proceedings of the National Academy of Sciences} 118(17):e2015972118.
\bibitem[Tkachenko et~al.(2021b)]{tkachenko2021elife} Tkachenko, A. V., Maslov, S., Wang, T., Elbanna, A., Wong, G. N. and Goldenfeld, N. (2021b). Stochastic social behavior coupled to COVID-19 dynamics leads to waves, plateaus, and an endemic state. \emph{eLife} 10:e68341.
\bibitem[Wallinga and Lipsitch(2007)]{wallinga2007} Wallinga, J. and Lipsitch, M. (2007). How generation intervals shape the relationship between growth rates and reproductive numbers. \emph{Proceedings of the Royal Society B} 274(1609):599--604.
\bibitem[Ward et~al.(2021)]{ward2021} Ward, H., Atchison, C., Whitaker, M. et al. (2021). SARS-CoV-2 antibody prevalence in England following the first peak of the pandemic. \emph{Nature Communications} 12:905.
\bibitem[Wong et~al.(2023)]{wong2023} Wong, K. L. M., Gimma, A., Coletti, P. et al. (2023). Social contact patterns during the COVID-19 pandemic in 21 European countries: evidence from a two-year study. \emph{BMC Infectious Diseases} 23:268.
\end{thebibliography}
\end{document}
